


\documentclass[11pt]{article}
\usepackage{amsgen,amsmath,amstext,amsbsy,amsopn,amsfonts,amsthm,amscd}
%\usepackage{bbm}
\usepackage[dvips]{graphicx}
\usepackage{enumerate}
\setlength{\textwidth}{16.0cm} \setlength{\textheight}{25cm}
\voffset=-2.50cm \hoffset=-2.00cm \setlength{\topmargin}{0.4cm}

%%%%%%%%%%%%%%%%%%%%%%%%%%%%%%%%%%%%%%%%%%%%%%%%%%%%%%%%%%%%%%

\newtheorem{lem1}{Lema}%[chapter]
\newtheorem{th1}{Teorema}%[chapter]
\newtheorem{prop1}{Proposici\'{o}n}%[chapter]
\newtheorem{defi1}{Definici\'{o}n}%[chapter]
\newtheorem{coro1}{Corolario}%[chapter]
\newtheorem{ob1}{Observaci\'{o}n}%[chapter]
\newtheorem{obs1}{Observaci\'{o}n}%[chapter]
\newtheorem{nota1}{Nota}%[chapter]
\newtheorem{eje1}{Ejemplo}%[chapter]
\newtheorem{ejes1}{Ejemplo}%[chapter]
\newtheorem{ejerc1}{Ejercicio}%[chapter]
\newtheorem{ejercs1}{Ejercicio}%[chapter]
\newenvironment{prop}{\begin{prop1}\rm}{\end{prop1}}
\newenvironment{defi}{\begin{defi1}}{\end{defi1}}
\newenvironment{coro}{\begin{coro1}}{\end{coro1}}
\newenvironment{th2}{\begin{th1}}{\end{th1}}
\newenvironment{nota}{\begin{nota1}\rm}{\end{nota1}}
\newenvironment{lem}{\begin{lem1}}{\end{lem1}}
\newenvironment{ob}{\begin{ob1}\rm}{\end{ob1}}
\newenvironment{obs}{\begin{obs1}\rm}{\end{obs1}}
\newenvironment{eje}{\begin{eje1}\rm}{\end{eje1}}
\newenvironment{ejes}{\begin{ejes1}\rm}{\end{ejes1}}
\newenvironment{ejerc}{\begin{ejerc1}\rm}{\end{ejerc1}}
\newenvironment{ejercs}{\begin{ejercs1}\rm}{\end{ejercs1}}
%%%%%%%%%%%%%%%%%%%%%%%%%%%%%%%%%%%%%%%%%%%%%%%%%%%%%%
%\DeclareSymbolFont{cucu}{U}{bbmss}{m}{n}
\DeclareMathOperator{\CC}{\mathbb C}%{\mathalpha}{cucu}{67}
\DeclareMathOperator{\NN}{\mathbb N}%{\mathalpha}{cucu}{78}
\DeclareMathOperator{\RR}{\mathbb R}%{\mathalpha}{cucu}{82}
\DeclareMathOperator{\QQ}{\mathbb Q}%{\mathalpha}{cucu}{81}
\DeclareMathOperator{\ZZ}{\mathbb Z}%{\mathalpha}{cucu}{90}
%%%%%%%%%%%%%%%%%%%%%%%%%%%%%%%%%%%%%%%%%%%%%%%%%%%%%%
\DeclareMathOperator{\sen}{sen}
\DeclareMathOperator{\arcsen}{arcsen}
\DeclareMathOperator{\senh}{senh} \DeclareMathOperator{\fr}{Frac}
%%%%%%%%%%%%%%%%%%%%%%%%%%%%%%%%%%%%%%%%%%%%%%%%%%%%
\newcommand{\ran}[1]{{\em Im\/}(#1)}
\newcommand{\ca}[1]{{\em card\/}(#1)}
\newcommand{\sop}[1]{{\em supp\/}(#1)}
\newcommand{\tra}[1]{{\em traza\/}(#1)}
\newcommand{\dis}[1]{d (#1)}
\newcommand{\inte}[1]{{\em Int\/}(#1)}
\newcommand{\conv}[1]{{\em conv\/}(#1)}
\newcommand{\too}{\longrightarrow}
\newcommand{\oo}{\overline}
\newcommand{\cq}{\begin{flushright} $\Box$ \end{flushright}}
\newcommand{\cqd}{\, \, \rule{7pt}{7pt}}
\def\limsup{\mathop{\overline{\rm lim}}}
\def\liminf{\mathop{\underline{\rm lim}}}
\renewcommand{\theenumii}{\arabic{enumii}}
\renewcommand{\labelenumii}{\theenumi.\theenumii.}
\renewcommand{\theenumiii}{\arabic{enumiii}}
\renewcommand{\labelenumiii}{\theenumi.\theenumii.\theenumiii}
%%%%%%%%%%%%%%%%%%%%%%%%%%%%%%%%%%%%%%%%%%%%%%%%%%%%%%%%%%%%%%%%%
\begin{document}

\section{COMPORTAMIENTO DE LOS SEGUNDOS AUTOVALORES MÁXIMOS.} 



\noindent Sea $\mathbf{M}$ una matriz hermitiana semidefinida positiva. 

\noindent Para cada $n\in \NN$ consideramos la sección de orden $n$, denotada por $\mathbf{M}_n$. 

\noindent Consideramos los autovalores de la sección $n\times n$ denotados de la siguiente forma 
 
$$
0 \leq \sigma_{n,n}\leq \dots \sigma_{2,n} \leq \sigma_{1,n}.
$$

\noindent Utilizando el teorema de entrelazamientos de Cauchy : 

\bigskip

\noindent {\bf Teorema de entrelazamiento de Cauchy} (\cite{Bhatia})Sea  $A$ una matriz de orden  $n$ y sea  $B$ una submatriz principal de  $A$ de  orden  $n-1$. Si $\lambda_n\leq \dots \leq \lambda_2\leq \lambda_1$ son los autovalores de $A$ y  $\mu_{n-1}\leq \dots \leq \mu_2\leq \mu_1$ son los autovalores de  $B$ entonces 
$$
\lambda_n \leq \mu_{n-1} \leq \dots \leq \mu_2 \leq \lambda_1.
$$


 \begin{prop} Sea $\mathbf{M}$ una matriz infinita  hermitiana semidefinida positiva. Sean 
 
$$
0 \leq \lambda_{n,n}\leq \dots \lambda_{2,n} \leq \lambda_{1,n}.
$$
\noindent los autovalores de la matriz $\mathbf{M}_n$. Entonces La sucesión $\{ \sigma_{2,n}\}_{n=1}^{\infty}$ es una sucesión creciente, es decir, 
$$
\lambda_{2,n} \leq \lambda_{2,n+1} \qquad \forall n\in \NN.
$$ 
\noindent por lo que existe $\displaystyle{\lim_{n\to \infty} \lambda_{2,n}}$ pudiendo ser finito o infinito. 
\end{prop} 


\begin{proof} Utilizando el teorema de entrelazamiento se tiene que: 

$$
0 \leq \lambda_{n+1,n+1} \leq \lambda_{n,n} \leq\lambda_{n,n+1} \leq \dots \lambda_{2,n}  \leq \lambda_{2,n+1} \leq \lambda_{1,n} \leq \lambda_{1,n+1}
$$

\noindent En particular, 
$$
\lambda_{2,n} \leq \lambda_{2,n+1} \qquad \forall n\in \NN.
$$
\end{proof} 

\noindent 


\noindent Cuando consideramos la matriz $\mathbf{D}$ infinita de Hessemberg, siguiendo la notación habitual que hemos introducido se tiene que: 

\noindent Para cada $n\in \NN$ la matriz $\mathbf{D}_n^{*}\mathbf{D}_n$  y denotamos por 

$$
0 \leq \sigma_{n,n}\leq \dots \sigma_{2,n} \leq \sigma_{1,n}.
$$

\noindent los valores singulares de la matriz $\mathbf{D}_n$ 
$$
\sigma_{k,n}=\sqrt{\lambda_{k,n}}
$$

\noindent siendo $\lambda_{k,n}$ los autovalores de la matriz $\mathbf{D}_n^{*}\mathbf{D}_n$. 

\noindent Se tiene el siguiente resultado: 

%\noindent Para una matriz finita o infinita denotamos por 

\begin{lem} Sea $\mathbf{D}$ una matriz infinita de Hessenberg superior y $\mathbf{D}_n$ las secciones de orden $n$. Entonces se tiene que 

 $$
(\mathbf{D}_{n+1}^{*}\mathbf{D}_{n+1})_{n} = \mathbf{D}_n^{*}\mathbf{D}_n +\begin{pmatrix}
0 & \dots&0 & 0\\
 & \ddots& & \vdots\\
  & \dots& & 0\\
0 & \dots& 0 & d_{n+1,n}^2

\end{pmatrix}
$$
\end{lem} 

 \begin{proof} Expresando el producto $(\mathbf{D}_{n+1}^{*}\mathbf{D}_{n+1})$ adecuadamente y multiplicando por cajas, obtenemos el resultado: 
    \[ (\mathbf{D}_{n+1}^{*}\mathbf{D}_{n+1}) = 
\left (   \mbox{
\begin{tabular}{ccc|c}
& $$ &  & $0$ \\
$$ & $\mathbf{D}_{n}^{*}$ &  & $\vdots$ \\
&  &  &$d_{n+1,n}^*$ \\  \hline

$d_{1,n+1}^*$ & $\ldots$ & $d_{n,n+1}^*$  & $d_{n+1,n+1}^*$ \\
\end{tabular}
} \right ) \left (   \mbox{
\begin{tabular}{ccc|c}
& $$ &  & $d_{1,n+1}$ \\
$$ & $\mathbf{D}_{n}$ &  & $\vdots$ \\
&  &  &$d_{n,n+1}$ \\  \hline

$0$ & $\ldots$ & $d_{n+1,n}$  & $d_{n+1,n+1}$ \\
\end{tabular}
} \right ) 
\]

\[
\left (   \mbox{
\begin{tabular}{ccc|c}
& $$ &  & \\
$$ & $\mathbf{D}_{n}^{*}\mathbf{D}_{n}+\begin{pmatrix}
0 & \ldots&0 & 0\\
 & \ddots& & \vdots\\
  & \dots& & 0\\
0 & \dots& 0 & d_{n+1,n}^2

\end{pmatrix}$ &  & $X$ \\
&  &  &$$ \\  \hline

$$ & $Y$ & $$  & $Z$ \\
\end{tabular}
} \right ) 
\]   
 
\end{proof}

\begin{prop} Sea $\mathbf{D}$ una matriz infinita de Hessenberg, y $\mathbf{D}_n$ las secciones de orden $n$ entonces si $$
0 \leq \sigma_{n,n}\leq \dots \sigma_{2,n} \leq \sigma_{1,n}.
$$

\noindent son los valores singulares de la matriz $\mathbf{D}_n$ se tiene que 
la sucesión $\{\sigma_{2,n}\}_{n=1}^{\infty}$ es creciente y existe
$$
\lim_{n\to \infty} \sigma_{2,n} \leq \infty
$$


\end{prop} 

\begin{proof}
   Sean $\lambda_{2,n}$ los segundos autovalores de una   matriz hermitiana de orden n. Denotamos por $H_n$  la siguiente matriz de orden $n$:
   \[
   \mathbf{H}_n=\begin{pmatrix}
0 & \ldots&0 & 0\\
 & \ddots& & \vdots\\
  & \dots& & 0\\
0 & \dots& 0 & d_{n+1,n}^2

\end{pmatrix}
   \]
   Teniendo en cuenta que $H_n$ es hermitiana semidefinida positiva y utilizando las propiedades de monotonía de Weyl (\cite{Bhatia}):
  \[
  \lambda_{2,n}(\mathbf{D}_n^{*}\mathbf{D}_n) \leq  \lambda_{2,n}(\mathbf{D}_n^{*}\mathbf{D}_n +\mathbf{H}_n)=\lambda_{2,n}(
\mathbf{D}_{n+1}^{*}\mathbf{D}_{n+1})_{n}   \] 
Ahora por las propiedades de entrelazamiento
\[
 \lambda_{2,n}(
\mathbf{D}_{n+1}^{*}\mathbf{D}_{n+1})_{n} \leq \lambda_{2,n+1}(
\mathbf{D}_{n+1}^{*}\mathbf{D}_{n+1}), \] 
de donde 
\[
 \lambda_{2,n}(\mathbf{D}_n^{*}\mathbf{D}_n) \leq \lambda_{2,n+1}(
\mathbf{D}_{n+1}^{*}\mathbf{D}_{n+1}). \] Por tanto $\sigma_{2,n}\leq\sigma_{2,n+1}$  y se verifica el resultado. 
\end{proof}

\bigskip
 

\begin{thebibliography}{5}

\bibitem[Bhatia97]{Bhatia}Rajendra Bhatia, \textsl{Matrix Analysis}, Springer, New York (1997).
\bibitem[Chi78]{chihara}Theodore S. Chiara, \textsl{An Introduction to
Orthogonal Polynomials}, Gordon and Breach, New York (1978).
\bibitem[Sze75]{szego}Gabor Szeg\"{o}, \textsl{Orthogonal Polynomials},
American Mathematical Society, Coloquium Publications, Vol. 32, primera
edici\'{o}n (1939), cuarta edici\'{o}n (1975).
\bibitem[Tor87]{tesis}E. Torrano, \textsl{Interpretaci\'on matricial de los
polinomios ortogonales en el caso complejo} Tesis Doctoral, Universidad de Santander,(1987).
\bibitem[Esc12]{tesis}C. Escribano, \textsl{Medidas Autosemejantes en el Plano Complejo, Momentos y Matrices de 
 Hessenberg} Tesis Doctoral, Universidad Politécnica de Madrid,(2012).
 
\bibitem[EST06]{} Escribano C. Escribano, M. Sastre, E.  Torrano.\textsl{ Moment matrix of self-similar measures.} Electronic Transactions on Numerical Analysis. Volume. 24. 79-87. (2006)
\bibitem[EGST11]{}
Escribano C. Escribano, A. Giraldo, M. Sastre, E.  Torrano. \textsl{Hessenberg matrix for sums of Hermitian positive definite matrices and
weighted shifts.}236(1), 98-106 (2011)
\end{thebibliography}

\end{document} 



\noindent INCLUIR CUANDO APARECE LA MATRIZ $\mathbf{D}$ y los valores singulares: 


\noindent Utilizando resultados conocidos se tiene que 

 $$
\Vert \mathcal{D}\Vert = \lim_{n\to \infty} \Vert \mathbf{D}_n \Vert = \lim_{n\to \infty} \sigma_{1,n} 
$$


\end{document} 




\section{UNA COTA PARA LOS CEROS DE LOS POLINOMIOS DE SOBOLEV} 

\noindent Producto Sobolev continuo: caso secuencialmente dominado. 

\noindent Estimaci\'{o}n de la cota de acotaci\'{o}n de ceros de polinomios ortogonales de Sobolev. 

\noindent Consideremos el siguiente producto de Sobolev: 
$$
<p(t),q(t)>=  \int_{a}^{b} p(t)q(t)w_0(t)\; dt +  \int_{a}^{b} p'(t)q'(t)w_1(t)\; dt
$$

\noindent donde $w_0(t),w_1(t)$ son pesos continuos de modo que est\'{a}n sequencialmente dominados en el sentido de que: Existe una constante $C>0$ tal que: 
$$
w_1(t) \leq Cw_0(t) , \qquad t\in [a,b]
$$

\noindent Puesto que $\dfrac{w_1(t)}{w_0(t)}$ es una funci\'{o}n continua en $[a,b]$ consideremos 
$$
\alpha= \max \{ \dfrac{w_1(t)}{w_0(t)}, t\in [a,b]\}.
$$

\noindent Consideremos el operador multiplicaci\'{o}n por $t$ 
$$
\mathcal{D}: \mathbb{P}[t]\to \mathbb{P}[t], \quad \mathcal{D}(p(t))=tp(t).
$$

\noindent Veamos que $\mathcal{D}$ es acotado en este caso adem\'{a}s: 
$$
 \Vert \mathcal{D} \Vert \leq C= (R^{2}+2\alpha)^{1/2}.
 $$

$$
\Vert \mathcal{D}(p(t))\Vert^{2} = \int_{a}^{b} t^2p^2(t)w_0(t)\; dt +  \int_{a}^{b} (tp'(t))^2w_1(t)\; dt= \int_{a}^{b} t^2p^2(t)w_0(t)\; dt +  \int_{a}^{b} (p(t)+tp'(t))^2w_1(t)\; dt \leq 
$$
\noindent utilizando que $(a+b)^2\leq 2(a^2+b^2)$ si $a,b\geq 0$ se tiene que: 
$$
 \int_{a}^{b} t^2p^2(t)w_0(t)\; dt+2\int_{a}^{b} p(t)^2w_1(t)\; dt +\int_{a}^{b} +t^2p'(t))^2w_1(t)\; dt \leq
 $$
 $$
  R^{2}\left( \int_{a}^{b} p^2(t)w_0(t)\; dt +  \int_{a}^{b} (p'(t)^2w_1(t)\; dt+ 2\int_{a}^{b} p(t)^2w_1(t)\; dt \leq 
$$
$$
R^{2} \Vert p(t)\Vert^{2} + 2\alpha \int_{a}^{b} p(t)^2w_0(t)\; dt \leq (R^{2}+2\alpha)\Vert p(t) \Vert^{2}
$$

\noindent Por ejemplo, si $[a,b]=[-1,1]$ y $\alpha =1$ quedar\'{\i}a $C=\sqrt{1+2}=\sqrt{3}$. 

\begin{enumerate} 
\item Experimentar con ejemplos sequencialmente dominados $w_0(t)=t^n$ y $w_1(t)=t^{2n}$ con $n$ par en el intervalo $[-1,1]$, se aprecia que la c ota es buena ???
\item Experimentar con ejemplos que no sean sequencialmente dominados $w_0(t)=t^{2n}$ y $w_1(t)=t^{n}$ con $n$ par en el intervalo $[-1,1]$, se aprecian resultados an\'{a}logos al caso anterior. 
\item Observar que en los casos anteriores todos los ceros están en el círculo cerrado de centro el origen y radio $\sqrt{3}$. 
\end{enumerate} 
\noindent 


\section{Conjetura (reunión 13 de marzo )} 

\noindent A la vista de los resultados: consideremos el conjunto de los ceros de los polinomios del polinomio ortonormal (ortogonal, etc.) de orden $n$ de Sobolev: 
$$
Z_n=\{ z\in \CC: P_n(z)=0\}= \{ z \in \CC: \vert \mathbf{D}_n - zI_n\vert=0\}.
$$

\noident es claro que el conjunto de todos los ceros de los polinomios de Sobolev es
$$
Z=\cup_{n=1}^{\infty} Z_n
$$

\noindent Por otra parte, si consideramos todos los autovalores de $\mathbf{D}_n$ son los ceros de $P_n(z)$, por lo que denotando el autovalor m\'{a}ximo de $\mathbf{D}_n$ como $\lambda_{max}(\mathbf{D_n})$ y denotando por $z_n$ el cero de $P_n(z)$ tal que 
$$
\vert z_n\vert =\lambda_{max}(\mathbf{D}_n)
$$

\noindent El problema de la acotación de los ceros consiste en el problema de la acotación de los $\{z_n\}_{n=1}^{\infty}$:



\bigskip

\noindent {\bf CONJETURA:} productos de Sobolev soportados en intervalos del tipo 
$[a_1,a_2],[a_3,a_4]$ si $R=\max \{\vert a_k \vert : k=1,\dots,4\}$
 
\begin{enumerate} 
\item 
$$
\lim_{n\to \infty} \vert \lambda_{max}(\mathbf{D}_n)\vert = R;
$$
 \item Si $\vert a_{k_0} \vert=R$ entonces 
 $$
 \lim_{n\to \infty} z_n = a_{k_0},$
 
\noindent es decir, $a_{k_0}$ sería un atractor de los ceros con módulo máximo. 

 \end{enumerate} 
 
\noindent Nótese que son equivalentes: 
\begin{enumerate} 
\item $Z$ es acotado.
\item La sucesión $(\lambda_{max}(\mathbf{D}_n))_{n=1}^{\infty}$ está acotada
\end{enumerate} 

\noindent Por lo tanto, se propone para el estudio de la acotación de los ceros de los polinomios ortogonales de Sobolev la acotación de estos autovalores máximos. 

\noindent En el caso de que las matrices $\mathbf{D}_n$ sean diagonalizables, lo que ocurre cuando \underline{todas} las raíces de los polinomios ortogonales son distintos. Esto ocurre en los ejemplos estudiados, pero sabemos que no es cierto en el caso de polinomios ortogonales de Sobolev. Tampoco es cierto en el caso de polinomios soportados en medidas en el plano complejo, pero sí lo es en el caso de medidas soportadas enla recta real, por la simetría de la matriz de Jacobi. 


\noindent No es cierto, en general, que 
$$
\Vert \mathbf{D}_n \Vert = \lambda_{max}(\mathbf{D}_n)
$$


\noindent Para ello, basta considera las secciones del shift-right. Por  ejemplo, 
la sección de orden dos, 
$$
\mathbf{D}_2 = \begin{pmatrix} 0 & 0 \\ 1 & 0 \end{pmatrix} 
$$

\begin{enumerate} 
\item $\Vert \mathbf{D}_2 \Vert =1 $ 
\item $\lambda_{max}(\mathbf{D}_2)=0$ (el único autovalor es $0$). 
\end{enumerate} 



\noindent Sin embargo, en el caso de que todas las raíces de los polinomios ortogonales de Sobolev (caso recta real) fuesen distintos entonces las matrices $\mathbf{D}_n$ serían diagonalizables y en este caso sí sería cierto que 
$$
\Vert \mathbf{D}_n \Vert = \lambda_{max}(\mathbf{D}_n)
$$
 
\noindent Se puede probar que para una matriz con entradas reales:
$$
\Vert \mathbf{D}_n \Vert= \lambda_{max}(\mathbf{D}_n^{t}\mathbf{D}_n)
$$

\noindent la matriz $\mathbf{D}_n^{t}\mathbf{D}_n$ es una matriz simétrica, autoadjunta, por lo que todos sus autovalores son reales. 
\bigskip

\noindent {\bf PROBLEMA:} ¿Es cierto para los casos de productos de Sobolev soportados en la recta real que los ceros de los polinomios ortogonales son todos distintos?  
\newpage

\section{Reunión de 27 de marzo} 


\noindent Comprobación de todos los pasos: 

\begin{prop} Sean $\mu_1,\mu_2$ medidas con soporte infinito  en $\CC$  y consideremos un producto interno 
$$
<p(z),q(z)>= \int p(z)\overline{q(z)} d\mu_0+ \int p(z)\overline{q(z)} d\mu_1
$$

\noindent Sea $\{P_n(z)\}_{n=0}^{\infty}$ la sucesión de polinomios ortonormales. Entonces 


$$
Z=
 \{ z \in \CC: \text{ there exists } n\in \NN, \varphi_n(z)=0\}\subset \{ z\in \CC  : \;  \vert z \vert < \Vert \mathcal{D} \Vert \}.
$$
\end{prop}

\begin{proof} Supongamos que  $z_0\in Z$, es una raíz de un polinomio $P_n(z)=(z-z_0)q_{n-1}(z)$; además   $q_{n-1} \neq 0$ y $q_{n-1}(z)$ es de grado menor o igual que $n-1$. 
$$
<P_n(z),P_n(z)>=<(z-z_0)q_{n-1}(z),(z-z_0)q_{n-1}(z)>=<zq_{n-1}(z),
$$



$$
\vert z_0 \Vert^2 \Vert q_{n-1} \Vert^2 = \Vert zP_{n-1} \Vert^2 - \Vert P_n(z) \Vert^2 < \Vert z P_{n-1}(z) \Vert^2 = \Vert \mathcal{D}(P_{n-1}(z)) \Vert \leq \Vert \mathcal{D} \Vert^2 \Vert P_{n-1} \Vert^2
$$

\noindent Since  $\Vert P_{n-1} \Vert \neq 0$ it follows that 
$$
\vert z_0  \vert^2 <  \Vert \mathcal{D} \Vert^2
$$

\noindent as we required.
\end{proof} 

\newpage 
\section{Reunión 10 de abril} 
\noindent Estudio de los casos 

$$
\Vert p \Vert^2_{\mathbf{S}} = \int_{-1}^{1} p^{2}(t)+ \int_{-1}^{1} t^{2} (p'(t))^{2}\; dt + \vert p'(0)\vert^{2}
$$

\noindent En el artículo J.Rodríguez, ...se afirma que el operador multiplicación por $z$ no es acotado, por lo que 
$$
\lim_{n\to \infty} \Vert \mathbf{D}_n \Vert =\infty
$$
 

\noindent En los siguientes casos, salen por Rodriguez??, también se puede afirmar lo mismo.  

$$
\Vert p \Vert_{\mathbf{S}} = \int_{-1}^{1} p^{2}(t)\; dt + \vert p'(a)\vert^{2}
$$
$$
\Vert p \Vert_{\mathbf{S}} = \int_{0}^{1} p^{2}(t)\; dt + \vert p'(a)\vert^{2}
$$

\noindent $\bullet$ {\bf ESTUDIAR EL COMPORTAMIENTO DEPENDIENDO DE LA LOCALIZACIÓN DEL ÁTOMO  $a$ : } 
$$
\lim_{n\to \infty} \rho(\mathbf{D}_n),
$$

\noindent Experimentalmente se observa que en el caso : 
$$
\int_{0}^{1} p^{2}(t)\; dt + \vert p'(0) \vert^{2},
$$
$$
\lim_{n \to \infty} \rho(\mathbf{D}_n)=0.
$$

\noindent {\bf INTENTAR PROBARLO} 



\noindent En algunos casos parece ir al átomo, en otros casos, al punto más alejado del origen?? La simetría parece tener que ver. 

\bigskip
\bigskip

\noindent $\bullet$ {\bf MUY INTERESANTE:}

\noindent Estudio de la convergencia del \underline{segundo} mayor   de los valores singulares de $\mathbf{D}_n$, que parece no ser el mismo que el del primer autovalor!!!!. 

\noindent Dada la matriz $\mathbf{D}_n^{t}\mathbf{D}_n$ que es simétrica y semidefinida positiva consideramos el segundo mayor valor singular, es decir, si ordenamos los valores singulares de esta forma:

$$
0\leq \sigma_{n,n} \leq \dots \leq \sigma_{2,n} \leq \sigma_{1,n} 
$$



\noindent son los valores singulares de $\mathbf{D}_n$, 

\bigskip

\noindent Nótese que $\Vert \mathbf{D}_n \Vert= \sigma_{1,n}$, ya que para una matriz simétrica y semidefinida positiva el mayor autovalor es la norma.  

\bigskip Por otra parte, es claro que 
$$
\lim_{n\to \infty} \sigma_{1,n}(\mathbf{D}_n)= \lim_{n \to \infty} \Vert \mathbf{D}_n\Vert
$$

\noindent que siempre existe pudiendo ser $\infty$ en el caso no acotado. 


\noindent $\bullet$ {\bf  ESTUDIO DEL COMPORTAMIENTO ASINTÓTICO de: } 
$$
\lim_{n\to \infty} \sigma_{2,n},
$$

\noindent tiene relación con el soporte de las medidas???
  
\noindent En los experimentos realizados se tiene que este límite existe. ¿Se puede justificar utilizando que las matrices $\mathbf{D}_n$ son secciones de una matriz que este límite existe? YO CREO QUE SÍ. 


\noindent El estudio del segundo valor singular de una matriz ( y de cualquiera de ellos) tiene una expresión min max por el Teorema de Courant Fisher

\noindent A partir de estos teoremas se obtiene la siguiente caracterización de dicho valor singular como el mínimo  de las normas del operador restringido a subespacios vectoriales de dimensión $n-1$ o codimensión $1$, es decir, 
$$
\sigma_{2,n} = \min_{S: dim(S)=n-1} \max_{x\in S,\Vert x\Vert=1} <\mathbf{D}_n^t\mathbf{D}_n x,x>= \min_{S: dim(S)=n-1} \Vert \mathbf{D}_n/S\Vert.
$$

\noindent En los experimentos se observa una convergencia de estos valores singulares, por ello se trataría de probar que 
existe
$$
\lim_{n \to \infty} \sigma_{2,n}
$$

\noindent podría haber un min-max infinito dimensinal en el sentido de considerar el mínimo de las normas del operador multiplicación por $z$ entre todos los subespacios de codimensión $1$. 

$$
\min_{S: cod(S)=1} \Vert \mathcal{D}/S \Vert 
$$

\noindent Puede ocurrir que 

$$
\min_{S: cod(S)=1} \Vert \mathcal{D}/S \Vert < \Vert \mathcal{D} \Vert = \infty.
$$

\noindent Sí, puede ocurrir, basta considerar un caso como una matriz que sea diagonal salvo en la primera fila que esté llena de unos, entonces las secciones tendrán normas que van a infinito pero si consideramos las restricciones en las sucesiones con $x_0=0$ se tendría que está acotada, y el mínimo anterior sería todavía menor. 


$$
\mathbf{D}= \begin{pmatrix}
  1  & 1 & 1& \hdots & 1 & \hdots  \\
  1       & \frac{1}{2} & 0  & \hdots & 0 & \hdots 
   \\
  0  & \frac{1}{2} & \frac{1}{3}  & \vdots  & \hdots & \hdots  
   \\
  0  & 0 & 0   & \vdots  & \hdots & \hdots  \
\end{pmatrix}
$$


\noindent {\bf Problema:} ¿Es 
$$
\lim_{n\to \infty} \sigma_{2,n}= \min_{S: cod(S)=1} \Vert \mathcal{D}/S \Vert ????
$$
\subsection{Acotación de $\mathcal{D}$ en hiperplanos en espacios de  Sobolev con peso} 

\noindent Un resultado muy interesante. 

\begin{prop} Sea $w(t)$ un peso continuo en $[a,b]$ y consideremos el producto Sobolev con peso en $0$ dado por:
$$
\Vert p(t) \Vert^{2} = \int_{a}^{b} p^{2}(t)w(t)\; dt + p'(0)^{2}.
$$



\noindent Se tiene : 
\begin{enumerate} 
\item El operador multiplicación por $z$ no es acotado, por lo que: 
$$
\lim_{n\to \infty} \sigma_{1,n} = \infty.
$$
\item  El operador multiplicación por $z$ es acotado en el subespacio infinito dimensional de codimensión $1$, $\mathbb{P}_{0}[z]$, y además
$$
\lim_{n\to \infty} \sigma_{2,n}= \min_{S: cod(S)=1} \Vert \mathcal{D}/S \Vert \leq \max\{\vert a \vert, \vert b \vert \} 
$$
\end{enumerate} 
\end{prop} 

\begin{proof} Consideremos el subespacio vectorial de codimensión $1$, $\mathbb{P}_{0}[z]$, sea $p(z)\in \mathbb{P}[z]$ con $p(0)=0$, se tiene que: 
$$
\Vert \mathcal{D}(p(t))\Vert^{2} = \int_{a}^{b} t^{2}p^{2}(t)\; dt + (tp(t))'(0)=  \int_{a}^{b} t^{2}p^{2}(t)\; dt + (p(t)+tp'(t))(0)\leq 
$$
$$
\max \{a^2,b^2\} \int_{a}^{b} p^{2}(t)\; dt + p(0)^2= \max \{a^2,b^2\} \int_{a}^{b} p^{2}(t)\; dt \leq 
$$
$$
\max \{a^2,b^2\} ( \int_{a}^{b} p^{2}(t)\; dt+p'(0)^2)=\max \{a^2,b^2\} \Vert p(t)\Vert^{2},  
$$

\noindent por continuidad se tiene que $\mathcal{D}$ es acotado en $\mathbb{P}_0[z]$ siendo su norma menor o igual que $R=\max\{\vert a \vert, \vert b \vert \}$. 
\end{proof} 

\noindent GENERALIZACIONES: 
\begin{prop} En un espacio de Sobolev con peso del tipo: 
$$
\Vert p(z) \Vert^{2} = \int \vert p(z) \vert^{2} d\mu +  \vert p'(0) \vert^{2}
$$

\noindent tal que $0$ no es punto de evaluación de la medida $\mu$  se tiene que: 

\begin{enumerate} 
\item $\mathcal{D}$ no es acotado, por lo que: 
$$
\lim_{n\to\infty} \sigma_{1,n}= \infty.
$$
\item $\mathcal{D}$ es acotado en algún hiperplano, por lo que: 

$$
\lim_{n\to \infty} \sigma_{2,n}= \min_{S: cod(S)=1} \Vert \mathcal{D}/S \Vert \leq \max\{\vert z \vert : z\in {\it supp}(\mu) \} 
$$
\end{enumerate}

\end{prop} 

\subsection{Estudio detallado del caso Legendre+peso en 0} 

\noindent Consideremos el siguiente producto de Sobolev: 
$$
\Vert p(t)\Vert^{2} = \int_{0}^{1} p^2(t)\; dt + p'(0)^{2} 
$$


\noindent Nuestro objetivo es probar de forma sencilla que: 
\begin{enumerate} 
\item $\mathcal{D}$ no es acotado y de hecho 
$$
\lim_{n \to \infty} \sigma_{1,n}=\infty
$$

\item 
 $$\min_{S: cod(S)=1} \Vert \mathcal{D}/S \Vert < \infty
 $$

\item $$¿\lim_{n\to\infty} \rho(\mathbf{D}_n)= ???$$
\end{enumerate}


\section{Problemas relacionados con el segundo valor singular máximo.}

\noindent Problemas:  Si tenemos una matriz hermitiana definida positiva, acotada, y consideramos las secciones, y sus valores singulares ordenados:
$$
0 \leq \sigma_{n,n}\leq \dots \sigma_{2,n} \leq \sigma_{1,n} 
$$

\begin{enumerate} 
\item ¿ es cierto que $\lim_{n\to \infty} \sigma_{2,n}$ existe??? 

\noindent La respuesta es positiva y es consecuencia del Teorema de Cauchy de entralazamiento de los autovalores. 

\noindent {\it Cauchy interlace theorem} Let $A$ be an Hermitian matrix of order $n$ and let $B$ a principal submatrix of $A$ or order $n-1$. If $\lambda_n\leq \dots \leq \lambda_2\leq \lambda_1$ lists the eigenvalues of $A$ and $\mu_{n-1}\leq \dots \leq \mu_2\leq \mu_1$ the eigenvalues of $B$ then 
$$
\lambda_n \leq \mu_{n-1} \leq \dots \leq \mu_2 \leq \lambda_1
$$

\noindent Este resultado justifica el entrelazamiento de las raíces de los polinomios ortogonales cuando consideramos medidas en el caso real puesto que las  matrices de Jacobi, que son las matrices de Hessemberg en estos casos, son hermitianas, y sus autovalores son las raíces de los polinomios ortogonales. 

\noindent Ahora puesto que $\mathbf{D}_{n-1}
$ es una submatriz principal de $\mathbf{D}_n$ con nuestra notación tenemos que: 
$$
\sigma_{n,n} \leq \sigma_{n-1,n-1} \leq \dots \leq \sigma_{2,n-1} \leq   \sigma_{2,n} \leq \sigma_{1,n-1} \leq \sigma_{1,n}.
$$

\noindent En particular, la sucesión $\{\sigma_{2,n}\}_{n=1}^{\infty}$ es una sucesión creciente, por lo que o converge o tiene límite infinito. EXISTE 
$$
\lim_{n\to \infty} \sigma_{2,n} \leq \infty.
$$
\item ?????Coincide con 
$$
\min_{V, cod(V)=1} \max \Vert \mathcal{D}/V \Vert 
$$
\noindent 

\item 
\end{enumerate} 

\noindent Sería interesante probar (o buscar referencias) de que 
$$
\lim_{n\to \infty} \sigma_{2,n} = \inf_{V, cod(V)=1} \sup \Vert \mathcal{D}/V \Vert 
$$


\newpage 

\section{CASO DE CIRCUNFERENCIAS} 



\noindent Vamos a considerar productos Sobolev en los que están involucrados medidas de Lebesgue. Denotamos ${\bf m}$ como la medida de Lebesgue normalizada en la circunferencia unidad donde: 
$$
<p(z),q(z)>=\int_{\mathbb{T}}p(z)\overline{q(z)}\; dz = 
\dfrac{1}{2\pi} \int_{0}^{2\pi} p(e^{it})\overline{q(e^{it}}\; dt,
$$

\noindent Utilizando 
$$
<p,q>=\int p(z) \overline{q(z)}d{\bf m} + \int_{\mathbf{S}_r} p'(z) \overline{q'(z)}d{\bf m_{0;r}
$$

$$
\Vert p(z) \Vert^{2}_{\mathbf{S}} = \Vert zp(z) \Vert_{{\bf m}}^{2} + \Vert p'(z) \Vert^{2}_{{\bf m}_{0;r}}=
$$






$$
\Vert \mathcal{D}(p(z))\Vert_{\mathbf{S}}^{2} = \Vert \vert zp(z)\Vert_{{\bf m}}^{2} + \Vert (zp(z))' \Vert^{2}_{{\bf m}_{0;r}}=
$$

\noindent Puesto que en la primera integral $\vert z \vert=1$ se tiene que: 

$$
=\Vert  p(z) \Vert_{{\bf m}}
^{2} + \Vert p(z)+zp'(z) \Vert^{2}_{{\bf m}_{0;r}}
$$
$$
\leq \Vert  p(z) \Vert^2_{{\bf m}} + 2 \Vert  p(z) \Vert^2_{{\bf m}_{0;r}} + 2\Vert zp'(z) \Vert^{2}_{{\bf m}_{0;r}} \leq \Vert  p(z) \Vert^2_{{\bf m}} + 2 \Vert  p(z) \Vert^2_{{\bf m}_{0;r}} + 2r^2\Vert p'(z) \Vert^{2}_{{\bf m}_{0;r}} \leq 
$$
\noindent Si $C= \max \{1,2r^2\}$ se tiene que 
$$
\leq C(\Vert  p(z) \Vert^2_{{\bf m}}+\Vert p'(z) \Vert^{2}_{{\bf m}_{0;r}})+2\Vert  p(z) \Vert^2_{{\bf m}_{0;r}}= C\Vert p(z) \Vert^{2}_{\mathbf{S}}+ \Vert  p(z) \Vert^2_{{\bf m}_{0;r}} 
$$

\noindent Utilizaremos ahora la desigualdad polinomial de Wirtinger para la medida ${\bf m}_{0;r}$ que dice que si $q(z)$ es un polinomio con $q(0)=0$ se tiene que: 
$$
\Vert q(z) \Vert^{2} = \int_{\mathbf{S}_r} \left(\sum_{k=0}^{n} a_kz^k \right)= \sum_{k=0}^{n} r^{2k} \vert a_k \vert^{2} \leq 
$$
$$
\sum_{k=0}^{n} r^{2k}  \vert a_k \vert^{2}\leq \sum_{k=1}^{n} kr^{2k-2}\vert a_k \vert^{2} = \int_{\mathbf{S}_r} \vert q'(z)\vert^{2}d{\bf m}_{0;r}=  \Vert q'(z)\Vert^{2}.
$$


\noindent Luego, 
$$
\vert q(0) \vert 
$$


\end{document} 
